\section{Limitations, Risks and History}
\label{sec:history}

\subsection{Risks}

\begin{description}
  \item[R1: sensitive paths are exposed by default.]  The root is the home folder, so without
        protection \code{\~{}/.ssh} and similar would be reachable.  Mitigated by the core deny
        rules, active by default and built in, and by a startup notice and a banner if the user
        weakens them.  Residual: a secret with an unremarkable name outside the core list, the optional
        rules left off, and a user who disables the core rules and ignores the warning.
  \item[R2: files with several names.]  Path rules cannot see a second name for a file.  Covered
        for denied names under the roots and for the credential locations, by refusing any file whose
        names are not all accounted for; see the limits below.
  \item[R3: untrusted content in a page that holds a session.]  Markdown, PDF and archive names are
        rendered or shown; mitigated by the layers of Section~\ref{sec:content}.
  \item[R5: another program reads the user's documents.]  Quick Look draws iWork and Office
        documents with Apple's generators.  Mitigated by giving it only a private copy of what the
        guard would serve, with bounds on size and time; residual: a flaw in a generator, as when
        Finder shows the same picture.
  \item[R4: permission prompts.]  macOS asks permission the first time for Documents, Desktop,
        Downloads and external volumes; the documentation says so.
\end{description}

\subsection{Known limitations}

\begin{itemize}
  \item A denied name outside every root and every protected credential location is not found as another
        name of a file, so a file with such a name is refused (a pnpm store outside the root, or a
        clone made with \code{git clone --local}, can make ordinary files unavailable).  A listing or a
        search may show a name or a size up to one second after the file was renamed onto a denied name;
        opening the file never does.  Cancellation of copying and listing is cooperative: a read that the
        file system is itself holding up cannot be interrupted, and keeps a Quick Look slot until it
        returns.  Search completeness does not include cloud-only folders, which are skipped on purpose.
  \item Response timing is not equalized: a denied path can be answered slightly faster than a missing
        one.  The status and body are identical.
  \item Anything that runs as the user can read the same files.  \texttt{fsb} is not a sandbox against
        local malware.
  \item The metadata endpoint reads a symbolic link's text from the requested path after the file was
        opened; an attacker who could rewrite that link during the call could learn the text of a denied
        path (not its contents).  Requires a same-user attacker.
  \item Rule files are read at startup only.
  \item On macOS the launch URL is visible to other users for the moment \code{open} runs, and
        no connection-owner check exists there.  On Linux, a port forward the user runs counts as
        the user, and where the owner of a connection cannot be determined (no \code{/proc}, WSL,
        some sandboxes) the launch URL is refused, so \texttt{fsb} cannot be used there.
  \item \code{--stop} would take another program of the same user for \texttt{fsb} if that
        program's executable file were itself named \code{fsb}.
  \item Under emulation (an amd64 container on Apple silicon), \code{--status} and \code{--stop}
        cannot recognize \texttt{fsb}: QEMU makes every emulated process's \code{/proc/PID/exe}
        name itself.
  \item Not yet tried with real cloud-placeholder files or in Firefox; a real mouse drag of a column
        header in Safari is unverified.
  \item The browser test for Safari occasionally fails in a run of tests that need a focused window,
        and passes on rerun.
\end{itemize}

\subsection{History of defects found by the design's own checks}

Stating the laws of the algebraic specification, auditing every comparison of names and paths,
and using the running application, and running it on Linux, found forty-eight defects, all fixed and each covered by a test where
the underlying condition could be reproduced at all (Section~\ref{sec:untested} says where it could not).
In brief:

\begin{enumerate}
  \item Simple case folding let \code{.\ss h} open \code{.ssh} (a bypass of a deny rule).
  \item A plain file named like a folder-only rule was listed but could not be opened.
  \item A negated glob class could match across a path separator.
  \item A row order that was not transitive when a time could not be parsed.
  \item The filter and the search disagreed about the sharp s.
  \item A stale fallback request could show the wrong file's preview.
  \item Firmlink prefixes were matched by byte length (a bypass under a root of \code{/}).
  \item Root containment compared spellings (wrong on a case-sensitive volume).
  \item Breadcrumbs, selection and the core-rule warning compared names by spelling.
  \item The credential rules were anchored to one home folder.
  \item A hard link to a credential file was served.
  \item The session cookie followed the host, not the port.
  \item A hard link to any other denied file was served.
  \item Rule files that loaded without error and did nothing.
  \item Secrets under other names (\code{id\_rsa}, certificate stores, password databases).
  \item File names that disguised themselves with bidirectional characters.
  \item Opening a UNIX domain socket used an errno (\code{EOPNOTSUPP}, or \code{ENXIO} on
        Linux) that the guard's error map did not recognize, so it surfaced as an opaque
        500 instead of the same ``not found'' every other special file gets.  Found from a
        screenshot of the running application, not from any test.
  \item The preview pane's Details panel shared one error variable with the main preview
        and hid its own message whenever the main preview had already failed, so it could
        read ``Loading…'' forever even though its own request had already finished.  Found
        alongside the previous defect, in the same screenshot.
  \item Listing or searching a folder that is itself a cloud-only placeholder (a File
        Provider domain such as OneDrive or iCloud Drive can make a whole folder dataless,
        not only the files in it) went straight into reading its directory entries, which
        the provider's virtual file system answered with something the guard's error map
        did not recognize, surfacing as an internal error instead of the same cloud
        message a dataless \emph{file} already gets.  Found from a screenshot of a real
        OneDrive folder, not from any test: no fixture anywhere had a real cloud-only
        placeholder, only a synthetic flag on individual files.
  \item The folder listing's own error mapping, separate from the preview pane's, had no
        case for status 409, so once the previous defect was fixed at the guard level the
        interface still showed ``Server error (409).'' for a cloud-only folder instead of
        explaining it.  Found alongside the previous defect.
  \item While fixing the above, a latent defect in \code{Search}: a directory reopened
        from its queue that turned out no longer to be a directory (a rare race) left its
        descriptor open, because the code path that skipped it predates the fix and never
        closed what \code{Open} had just handed it.  The new dataless check would have
        turned this from a rare race into a routine leak, since a cloud-only subfolder
        hits the same code path on every walk that reaches one.
  \item Found by the user: a root given twice (\code{--root \~{}} when the home folder is
        already the default root) was served twice, and the interface, which keys its list of
        roots by path, failed so that every folder looked empty.  Roots are now served once,
        compared by identity.
  \item Found by the user: Chrome's PDF viewer moves the focus to the password field of an
        encrypted PDF, after which no key reached the interface, so arrowing through a folder
        stopped dead on that file.  Focus now returns to the list unless the user clicked or
        tabbed into the viewer.
  \item Found by the second independent review (and confirmed): a Finder alias named like an
        Office document made Quick Look draw the alias's target, which could be denied or outside
        the roots.  Quick Look now draws a private copy (Section~\ref{sec:quicklook}).
  \item Found by the same review: refusing a package document that held a denied part answered
        differently from a clean package, revealing what a listing hides.  Such parts are now
        left out of the copy.
  \item Found by the same review: \code{--stop} took any process whose \code{argv[0]} ended in
        \code{fsb} for \texttt{fsb}.  It now checks the owner, the executable file and the
        recorded start time (Section~\ref{sec:port}).
  \item Found by the same review: \texttt{fsb}'s own files in \code{\~{}/.config/fsb} were
        written through symbolic links.  They are now opened without following links.
  \item Found by the first run of continuous integration: the module named Go 1.26.0, whose
        standard library has since had security fixes (reported by \code{govulncheck}); the
        developer's machine happened to use a newer release.  The module now names the
        toolchain (Go 1.27.1).  The same run found two tests that assumed macOS.
  \item Found by an external review (October 2026, eleven findings, all confirmed and fixed;
        the commits cite them by number).  The ones that mattered most: a file with several names that
        the index had not seen was \emph{allowed} during the rebuild cooldown (the code returned the
        wrong flag, and the tests, which set the cooldown to zero, hid it); the index trusted names
        recorded at the last walk; a zip's entry limit compared a 16-bit count with 200{,}000, so it
        could never trip, and the standard library built the whole directory first; the index walk and
        the Quick Look package copy entered cloud-only folders; a symbolic link to \code{/} passed a
        check that only looked at spellings; an image fetch read its whole body before checking its
        size; search's work limit ignored entries it excluded; a tar.gz stopped exactly at a member
        boundary looked complete; an old search's timer overwrote a new search's results; resize-handle
        keys also moved the selection; and a collator that compares digits by value left
        \code{file1} and \code{file01} in an order that depended on the input.
  \item Found by an independent review of those fixes: two more hard-link bypasses and a quadratic
        cost in the index, and zips written to a pipe were refused.  Fixing them with a change-time test
        was itself broken by the next review (renaming a folder renames its files without touching
        them), which led to the policy of Section~\ref{sec:linkindex}: serve a file only when every name
        is accounted for, now.  A third round refined it (names folded by parent identity; one
        retry walk).
  \item Found by a second review of the whole (seven findings, all medium): a hard-link verdict
        remembered for the length of a listing or search hid a rename that kept the link count; a
        folder queued for search could be swapped for a symbolic link before it was entered; the
        one retry walk a short-counted file is owed was spent by a lookup during the cooldown that
        could not start it; the wrapper that added cancellation hid \code{Seek}, so listing a plain
        tar read every member's body; a copy that ran out of time while a read was blocked reported
        success when the read returned end of file; a folder that could not be opened left a search
        looking complete; and an empty extended attribute that grew between two system calls was
        sliced out of an empty buffer and panicked.  A third review confirmed each fix and found
        nothing new.  Two choices came out of it and are stated as limits rather than hidden: the
        one-second reuse of a verdict, and cooperative (not hard) cancellation.
  \item Found by the Linux distribution tests (October 2026; fixed in 1.0.2): the core rules
        named only the macOS locations of browser data, so on Linux the Chrome, Chromium and
        Firefox profiles, the GNOME and KDE keyrings and the \code{pass} store were served; and
        the identity protection of credential folders looked for other homes only in
        \code{/Users}, not \code{/home}.
  \item Found while fixing the previous one: Brave and Edge were protected on neither system
        (1.0.2); and, once the Linux rules were in, Flatpak keeps an application's settings in
        \code{\~{}/.var/app/ID/config}, which they did not reach (1.0.4).
  \item Found by the same tests: \code{--status}, \code{--stop} and the port-holder message ran
        \code{ps} and \code{lsof}, which minimal Debian, Fedora and Arch lack and BusyBox's
        \code{ps} cannot serve, so \code{--stop} said no \texttt{fsb} was running while one
        served; and a running \texttt{fsb} whose file had been replaced, as by a rebuild, was not
        recognized, because Linux appends \code{" (deleted)"} to \code{/proc/PID/exe}.  Linux now
        reads \code{/proc} (1.0.3).
  \item Found by the same tests: a name that is not UTF-8 reached the client as U+FFFD, under a
        name that does not exist, and could not be opened (1.0.3,
        Section~\ref{sec:wirenames}).
  \item Found by probing Quick Look on Linux: with no Quick Look, \texttt{fsb} still copied the
        document and then ran any program called \code{qlmanage} on \code{PATH} (1.0.5).
  \item Found by an independent review of findings 32 to 36 (October 2026; 1.0.6): the
        Linux rules still missed Snap Chromium (\code{\~{}/snap/chromium/common/chromium}),
        although this document and the README said the rules covered Snap; Firefox 147's new
        profile location (\code{\~{}/.config/mozilla}); Flatpak apps' own keyrings; and, on
        macOS, Chromium and Chrome Beta and Canary.  Vivaldi, Opera, Thunderbird and the NSS
        key store were added with them.
  \item Found by the same review: the PID file recorded the start time as local-time text, so
        \code{--stop} under another \code{TZ} did not find a background \texttt{fsb} and
        removed its PID file (on macOS since 1.0).  The record is now independent of the time
        zone (Section~\ref{sec:port}).  The same review found that one half of the test of
        \code{--stop}'s impostor check was refused by the start-time check before it reached
        the impostor check, and that the preview pane's Copy path skipped the shell quoting
        of finding 35.
  \item Found by the vulnerability scan of continuous integration on 2026-10-09: six advisories against the
        standard library's \code{net/http} in Go 1.27.1 (GO-2026-6603, -6609, -6611, -6612,
        -6613, -6617).  Four concern the HTTP/2 server, which \texttt{fsb} never uses; the
        unbounded \code{Range} header was reachable through downloads with a session.  Fixed by
        Go 1.27.2 (1.0.7).
  \item Found by installing 1.0.7 the way a user does: \code{go install} ignores a module's
        \code{toolchain} line, so the pinned Go 1.27.2 reached no one who installed a release, and
        every earlier release had been built with whatever Go the user had.  The module's
        minimum (\code{go} line) is now 1.27.2, which an older Go downloads (1.0.8).
  \item Found by an independent review of opening the browser on Linux (1.0.9): the background
        child passed its handoff pipe and its marker variable to the browser it started, so a
        browser holding the pipe kept \code{fsb --background} waiting, and an \texttt{fsb} later
        started from the browser's process tree took itself for a background child, printing
        nothing and taking over the PID file (Section~\ref{sec:port}).
  \item Found by the same review: without a display, \code{xdg-open} ran \code{w3m}, which fetched
        the launch URL itself and left the printed one used up, with no message.
        \code{xdg-open} is no longer run without a display.
  \item Found by the same review: the launch URL, in the command line of \code{xdg-open} and the
        browser, could be read by another user, who could use it first and take the session.
        First documented (1.0.9), then closed on Linux by accepting the launch only from a
        connection of \texttt{fsb}'s own user (Section~\ref{sec:launchowner}, 1.0.10).
  \item Found by the same review: Control-C and \code{fsb --stop} waited out the opener's grace
        period, and a background instance stopped during it was reported as running.
  \item Found by an independent review of the launch-owner check: after the user's browser had
        used the token, another user could present it without limit, each time making
        \texttt{fsb} read the connection table and print a warning (about 4\,GB of log an hour).
        A used token is now refused before any lookup, and warnings are limited.
  \item Found by the same review: the check took the overflow uid, which every user an unusual
        user namespace cannot name shares, for an owner; missed a client that closes its sending
        side after its request; reported an unreadable table as a missing row; and, with
        \texttt{fsb} run as root, blamed ``another user'' for the user's own browser.  The tests
        compared owners by uid only, so a lookup that matched \texttt{fsb}'s own end of the
        connection, letting anyone through, would have passed them; they now compare socket
        inodes.
  \item Found by an independent review of the whole core before the final release
        (reproduced with Chromium and two users): the launch URL named the path prefix,
        and since 1.0.9 it was in a command line other users can read.  A server of another
        user on another port, asked by a page the user's browser visited once for a URL
        under the prefix, received the session cookie and could replay it.  The launch URL
        now names no prefix, and on Linux a request with the cookie must come from the
        user's own connection (Section~\ref{sec:launchowner}, 1.0.11).
  \item Found by the same review: under \code{--root /}, \code{/.nofollow/}\emph{path}
        named a file a second time and kept the prefix through symbolic-link resolution,
        so a deny rule anchored to a place did not match it (Section~\ref{sec:aliases},
        1.0.11).  The differential test had tried every alias but only against secrets that
        rules matching a name anywhere protect.
\end{enumerate}

The lesson common to several: name equality is a property of the volume, so deny checks may over-match but
allow checks must decide by identity.  Findings 17, 18, 19 and 20 are a different lesson: an error path that
no test happened to reach (no fixture had ever contained a real socket, a real cloud-only folder, or a
selected file whose main preview fails) is still worth trying by hand.  Findings 29 to 31 are a fourth: a fix can introduce the next defect (a wrapper that hid
\code{Seek}, a retry spent before it ran, a memo that outlived the facts it summarized), so each
fix got a test that fails without it, and the reviews were repeated until one found nothing.  Finding 24 is a third: handing a
path to another program hands over that program's idea of what the path means, which the guard cannot
check; the safe input is bytes the guard has already read.  Findings 32 to 36 are a fifth: a test suite that
passes on a platform is not use of the program there.  The Go tests had run on Linux in continuous
integration since the first day, yet each of these appeared only when \texttt{fsb} was run on Linux
distributions, as a user would; the same work found that two tests had never tested what they
claimed there (the extended-attribute test skipped on the containers' file system, and the test of
\code{--stop}'s impostor check used \code{sleep} and \code{nc}, which on BusyBox exit at once when
started under another name).  Findings 37 and 38 add a sixth: a claim in the
documentation is as much in need of review as the code; ``covers Snap'' had been written,
and read, without anyone listing which Snap packages keep which folder.  Findings 39 to 48
add a seventh: what reaches the user is the release as installed, and the programs
\texttt{fsb} starts.  Every check had passed inside the repository, where the
\code{toolchain} line applies; only installing the tag as a user would showed that it
did not.  A program \texttt{fsb} starts inherits more than its arguments (descriptors,
environment, a terminal), and is itself a channel another user can read.

\subsection{Milestones}

\begin{center}
\small
\begin{tabular}{@{}p{0.16\linewidth}p{0.36\linewidth}p{0.36\linewidth}@{}}
\toprule
Milestone & Scope & Status \\
\midrule
M0 Skeleton & Repository, CI, guard, matcher, middleware & Done. \\
M1 Crawl & Browsing, listing, filtering, sorting, rules & Done; 100{,}000-file folder measured. Daily use by the author still to confirm. \\
M2 Walk & Preview, hover peek, search, attributes, columns, rendered Markdown, PDF, archives & Built and tested in Chrome and Safari; not yet on real cloud placeholders or in Firefox. \\
Hardening & The audits and defects of the preceding subsection & Done, including independent reviews in October 2026 of 1.0 (three rounds, the last finding nothing new), of the Linux work, of opening the browser on Linux, of the launch-owner check, and of the whole core before the final release, and continuous integration on macOS and Linux.  A professional penetration test was planned but, by the owner's decision, not done. \\
M3 Public release & Public repository, documentation, a source release & Done: version 1.0 (October 2026), source only; no binaries and no Homebrew formula, by decision. \\
Linux check & Five distributions in Docker; \code{/proc}, \code{/sys}, \code{/dev}; extended attributes; Quick Look; the browser suites against four distributions; an independent review & Done, October 2026: findings 32 to 38, fixed in 1.0.2 to 1.0.6. \\
Linux browser & Opening the browser with \code{xdg-open}; the launch only from the user's own connection; Go 1.27.2 & Done, October 2026: findings 39 to 48, fixed in 1.0.7 to 1.0.11. \\
M4 Run & The stage-three requirements of the functional specification & Not planned: the project is finished and provided as-is. \\
\bottomrule
\end{tabular}
\end{center}
